\documentclass[11pt,a4paper,roman]{moderncv}
\moderncvstyle{banking}
\moderncvcolor{purple}
\nopagenumbers{}
\usepackage[utf8]{inputenc}
\usepackage[T1]{fontenc}
\usepackage{lmodern}
\usepackage[scale=0.96,top=0.5cm,bottom=0.5cm,left=0.75cm,right=0.75cm]{geometry}
\renewcommand{\baselinestretch}{0.93}
\setlength{\hintscolumnwidth}{2.25cm}
\setlength{\footskip}{18pt}
\name{Wahid}{SAMER}
\address{Cergy (95), France}{}{}
\mobile{07 53 10 95 74}
\email{samer.ahmed.wahid@gmail.com}
\social[linkedin]{Wahid SAMER}
\social[github]{wahidrt}
\title{Alternance M2 - Quant Developer}
\begin{document}
\makecvtitle
\vspace{-3.3em}
\begin{center}
\parbox{0.97\textwidth}{\centering\small\itshape
Profil \textbf{mathématiques + développement}, orienté programmation scientifique, moteurs de backtest et implémentation de modèles quantitatifs en Python/C++. Recherche d'une \textbf{alternance M2 de 12 mois à partir du 21 septembre 2026}. Mobilité : \textbf{France, Luxembourg, Suisse, Belgique}.}
\end{center}
\vspace{0.15em}

\section{Formation}
\cventry{2025 -- 2027}{Master Mathématiques Appliquées à l'Ingénierie Financière}{CY Cergy Paris Université / CY Tech}{Cergy}{}{\textbf{M1 :} programmation Python, C++, Visual Basic, simulation stochastique, processus stochastiques, méthodes numériques avancées et optimisation.\\
\textbf{M2 2026--2027 :} model calibration and simulation, stochastic calculus, séries temporelles et machine learning avec Python.}
\cventry{2021 -- 2025}{Licence Mathématiques et Applications}{Université d'Évry Paris-Saclay}{Évry}{}{Formation solide en mathématiques pures et appliquées : probabilités, mesure et intégration, analyse, algèbre, topologie, analyse numérique et programmation.}

\section{Projets ciblés}
\cventry{2026}{Python / Git}{Architecture de trading quantitatif multi-actifs}{Projet M1 - équipe de 4}{}{Conception modulaire : chargement des données, modèles quantitatifs, portefeuille, simulation d'ordres, coûts, backtesting, métriques de performance et préparation d'une connexion MetaTrader 5.}
\cventry{2026}{Python}{Modèles de séries financières}{Projet M1}{}{Implémentation/étude de cointégration, VECM, filtre de Kalman, Ornstein--Uhlenbeck et changements de régime ; tests unitaires et échanges de données entre modules.}
\cventry{2024}{C / SDL2}{Moteur de jeu Tetris}{Projet Licence 3}{}{Programmation en C, logique temps réel, structures de données, gestion mémoire, événements et interface graphique SDL2.}

\section{Compétences}
\cvitem{Développement}{Python (NumPy, Pandas, SciPy, Scikit-learn, Matplotlib), C, C++, VBA, SQL, Julia.}
\cvitem{Ingénierie}{Git/GitHub, tests unitaires, architecture modulaire, traitement de données, backtesting, \LaTeX.}
\cvitem{Quant}{Processus stochastiques, séries temporelles, optimisation, méthodes numériques, Black--Scholes, Kalman, VECM.}
\cvitem{Langues}{Français : courant ; Arabe : langue maternelle ; Anglais : intermédiaire.}

\section{Expérience \& Engagement}
\cventry{Mai -- juin 2025}{Professeur de mathématiques - stagiaire}{Lycée Pierre Mendès France}{Vitrolles}{}{Préparation et transmission de notions de niveau Terminale ; développement de la rigueur, de la pédagogie et de la communication scientifique.}
\cventry{2019}{Technicien aéronautique - stagiaire}{G1 Aviation}{Gap}{}{Assemblage/maintenance dans un environnement exigeant en précision, procédures et sécurité.}
\cvitem{Associatif}{Bénévole aux \textbf{Restaurants du Cœur} pendant près de 3 ans : solidarité, accueil et travail d'équipe.}
\end{document}
